\begin{abstract}
Tool-using agents pay for their tools with context. Every Model Context Protocol (MCP) server an agent connects to contributes tool schemas up front and tool results as work proceeds, and general-purpose database servers can return arbitrarily large result sets. This report studies OraViz, a minimal, visualization-first MCP server for Oracle AI Database, and treats the context cost of a server's design choices as a measurable property. Three themes organize the discussion: a context contract that bounds every tool result through preview caps with explicit truncation markers, per-cell truncation, summarized large values, and profile-first chart selection; an architecture of seven single-purpose tools built on thin-mode \texttt{python-oracledb} and FastMCP; and an evaluation methodology in which identical questions are posed to OraViz and to the official Oracle SQLcl MCP server over the same Oracle AI Database 26ai Free instance, with tool schemas and tool results counted in \texttt{tiktoken} tokens. Across a five-question workflow, OraViz reduces the token budget by $53.3$\% overall: $59.7$\% fewer tokens for tool schemas, $80.5$\% for schema discovery, and $61.4$\% for an unaggregated $96$-row dump, while small result sets pay a bounded markdown framing overhead of roughly $55$\% that cannot grow with result size. Rendering the workflow aggregate as a chart costs $123$ text tokens plus one PNG. The report closes with reproduction-oriented appendix material covering container setup, demo data, and the benchmark command.
\end{abstract}
